\subsection{The affine abelian cover image retains exactly one cusp parameter}

\begin{statement}{A.167}{Proved}
{Complete classification in the affine abelian quotient of the degree-eight cover}
Let \(H=\pi_1(D^\circ_{\overline\Q})\subset\Pi\) be the
geometric subgroup of the actual degree-eight cover in A.139, and put
\[
 K_{\rm aff}=\overline{[H,H]},\qquad
 Q_{\rm aff}=\Pi/K_{\rm aff}.
\tag{A.167.1}
\]
This is an arithmetic normal kernel.  The geometric quotient has the
literal presentation
\[
 Q_{\rm aff}\simeq\widehat\Z[C_8]\rtimes\widehat\Z,
 \qquad x\longmapsto(e_0,2),\quad y\longmapsto(0,1),
\tag{A.167.2}
\]
where the second factor acts on the eight basis vectors by cyclic shift
through its reduction modulo eight.

Choose the rational local parameters \(W\) at
\(P_0,P_1,P_{1/2}\), and \(W/T\) at \(P_\infty\).
They fix tangential splittings at the four boundary points.  Write
\(B=\{P_0,P_1,P_{1/2},P_\infty\}\).  Relative
to these splittings, the images of full cusp sections give a bijection
\[
 \begin{gathered}
 \mathcal S_{\rm aff}(D^\circ)=
 \left\{\begin{array}{c}
 \text{sections of }E_D/K_{\rm aff}\text{ locally matching}\\
 \text{actual point or cusp sections}
 \end{array}\right\}\big/H^{\rm ab},\\
 \coprod_{P\in B}\widehat{\Q^\times}
           \xrightarrow{\ \sim\ }\mathcal S_{\rm aff}(D^\circ).
 \end{gathered}
\tag{A.167.3}
\]
The quotient on the right is conjugacy by the actual geometric
abelian group of the selected cover.  In particular both its boundary
point and its completed tangent parameter are determined by its class.

Consequently every locally point/cusp-matching section of
\(E_U/K_{\rm aff}\) in the combined modulo-eight class
\([64a^2]\) is geometrically conjugate to the image of an actual
rational cusp section of \(U\).  After choosing a global fixed sheet,
the sharper classification is (A.167.3).  Downstairs the exact set of
these geometric section classes is
\(\coprod_{P\in B}\widehat{\Q^\times}/\{\pm1\}\):
the four boundary labels remain distinct, and the additional
identification at each is precisely \(\tau\sim-\tau\).
For a full metabelian section
this is a conclusion about its image under the actual map
\[
             E_U/\Pi''\longrightarrow E_U/K_{\rm aff};
\tag{A.167.4}
\]
no equality with a cusp section in the larger source quotient is
asserted.
\end{statement}

\begin{proof}
\emph{The geometric quotient and the compactification map.}
Use the actual Nielsen basis \(u=xy^{-2},v=y\).  The character of
the cover sends \(u\) to zero and \(v\) to one in \(C_8\).
The Schreier basis of \(H\) is
\[
          u_j=v^juv^{-j}\ (0\leq j\leq7),\qquad z=v^8.
\tag{A.167.5}
\]
Its abelianization is \(\widehat\Z^9\) on these generators.
Conjugation by \(v\) shifts the \(u_j\) cyclically, since
\(v u_7v^{-1}=z u_0z^{-1}\), and fixes \(z\).
The homomorphism in (A.167.2) sends these generators to
\((e_j,0)\) and \((0,8)\).  Its image of \(H\) is the
abelian open subgroup
\(\widehat\Z[C_8]\times8\widehat\Z\), and its map on
\(H^{\rm ab}\) is an isomorphism: multiplication by eight identifies
\(\widehat\Z\) homeomorphically with \(8\widehat\Z\).
The two displayed images of \(x,y\) topologically generate the
target.  The resulting quotient map is therefore onto, and its kernel
inside the open subgroup \(H\) is exactly \(K_{\rm aff}\).
This proves the full topological isomorphism (A.167.2).

The cover is defined over \(\Q\) and is a geometric cyclic-eight
cover, so \(H\) is normal in \(E_U\).  Its closed commutator is
characteristic in \(H\), proving the arithmetic normality in
(A.167.1).  Since \(\Pi/H\) is abelian, \(\Pi'\subset H\)
and hence \(\Pi''\subset K_{\rm aff}\), giving (A.167.4).

Orient the four boundary meridians by the chosen local parameters.
In (A.167.5), their exact abelian vectors are
\[
 \begin{aligned}
 e_{P_0}&=\sum_{j\ {\rm even}}u_j+z,
 &e_{P_1}&=\sum_{j\ {\rm odd}}u_j+z,\\
 e_{P_{1/2}}&=z,
 &e_{P_\infty}&=-\sum_{j=0}^7u_j-3z.
 \end{aligned}
\tag{A.167.6}
\]
Indeed they are the words \(x^4,yx^4y^{-1},y^8,(xy)^{-8}\)
from the actual inertia orbits of the cover.  The last word uses
\(xy=uv^3\); its eighth power visits each index once and has final
factor \(v^{24}=z^3\).  The four vectors sum to zero and span a
primitive rank-three submodule.  For example their first three are a
basis of that submodule, since the even lamp sum, odd lamp sum and
\(z\) can be completed to an integral basis of \(\widehat\Z^9\).

Compactification kills precisely these four inertia groups.  The
induced abelian map is the exact sequence
\[
 \begin{gathered}
 0\longrightarrow M
   \xrightarrow{\ e_P\mapsto e_P\ }H^{\rm ab}
   \longrightarrow TJ\longrightarrow0,\\
 M=\frac{\displaystyle\bigoplus_{P\in B}\widehat\Z(1)e_P}
                 {\widehat\Z(1)\sum_{P\in B}e_P}.
 \end{gathered}
\tag{A.167.7}
\]
Here the right map is induced by the open immersion into the proper
curve and its Abel--Jacobi map, as in A.163.  Exactness on the left
can also be seen directly from the primitive vectors (A.167.6): the
kernel of the abelian compactification map is the closed span of the
boundary words, and its sole displayed relation is their sum.
All four points are rational.  Arithmetic conjugation takes an inertia
generator to a conjugate of its cyclotomic power; conjugation by
\(H\) disappears in \(H^{\rm ab}\).  Thus (A.167.7) is an exact
sequence of actual \(G_\Q\)-modules, with the diagonal cyclotomic
action on \(M\).  It is the toric boundary sequence, with its
inclusion and quotient maps specified.

\emph{Tate invariants and the continuous Kummer comparison.}
For \(k=\Q\), and for every finite completion \(k=\Q_p\),
\[
                              (TJ)^{G_k}=0.
\tag{A.167.8}
\]
For \(\Q\) the torsion is the finite group in A.164.  Locally a
sufficiently small formal subgroup of \(J(\Q_p)\) is, under the
logarithm, a torsion-free \(\Z_p\)-lattice of finite index.  A
torsion subgroup injects into the finite quotient by that subgroup and
is therefore finite.  A compatible Tate-module invariant would give
a compatible inverse system of points in this finite torsion group.
For any fixed level, transition from a sufficiently divisible higher
level multiplies by its exponent and is zero.  Every coordinate is
therefore zero, proving (A.167.8).  The exact cohomology sequence of
(A.167.7), or its direct cocycle proof, now gives injectivity
\[
          H^1_{\rm cts}(k,M)\longrightarrow
                         H^1_{\rm cts}(k,H^{\rm ab}).
\tag{A.167.9}
\]
In the direct proof, a coboundary in \(H^{\rm ab}\) whose values
lie in \(M\) comes from an element whose image in \(TJ\) is
invariant; (A.167.8) puts that element in \(M\) itself.

The Kummer identification used below has exactly the completed domain
\[
 H^1_{\rm cts}(k,\widehat\Z(1))
       \simeq\varprojlim_m k^\times/(k^\times)^m
       =\widehat{k^\times}.
\tag{A.167.10}
\]
At finite level this is the usual connecting map for multiplication
by \(m\) on \(\mathbf G_m\).  Compatible finite classes have
compatible continuous cocycle representatives: the set of
representatives of each fixed class is a finite zero-cochain orbit,
and a common multiple level solves any finite string of compatibility
conditions.  Compactness of these finite sets gives the compatible
system.  Conversely, if two continuous cocycles have equal finite
classes, the finite sets of conjugating zero-cochains have the same
finite-intersection argument.  This proves (A.167.10) and its
compatibility with every localization.  In particular no replacement
by rational elements, fixed finite support, or one finite exponent is
being made here.

\emph{Recovering the unique global inertia line.}
Let \(s\) be a section in the right side of (A.167.3).  Its proper
image satisfies A.163.  That theorem gives a single \(P\in B\)
such that the proper section class is the image at \(P\), and every
selected proper finite local centre is \(P\).  Fix the tangential
section \(\sigma_P\) associated with the specified rational local
parameter.  It splits \(E_D/K_{\rm aff}\), whose geometric group
is \(H^{\rm ab}\).  Conjugate \(s\) by a lift in
\(H^{\rm ab}\) of its proper conjugator in \(TJ\).  Its proper
image is then literally the proper image of \(\sigma_P\).
Writing their difference in this splitting gives a continuous
cocycle with values in the kernel of (A.167.7):
\[
                    s(\gamma)=c(\gamma)\sigma_P(\gamma),
                    \qquad c\in Z^1_{\rm cts}(\Q,M).
\tag{A.167.11}
\]
Its class in \(H^1(\Q,M)\) is determined by the original
section class because of (A.167.9).

At a finite place the actual local image is a cusp at this same
\(P\).  It cannot be an ordinary point of \(D^\circ\): its
compactified point would then differ from \(P\), contrary to
A.163 and local full Kummer injectivity.  A local cusp at \(P\)
differs from the fixed tangent section by the Kummer class of its
local tangent parameter in the actual inertia line
\(I_P^{\rm ab}=\widehat\Z(1)e_P\).  Consequently the class of
\(c_p\) in \(H^1(\Q_p,H^{\rm ab})\) is the image of a class
in \(H^1(\Q_p,\widehat\Z(1)e_P)\).  Injectivity (A.167.9)
makes the same equality hold in \(H^1(\Q_p,M)\).

The quotient \(M/\widehat\Z(1)e_P\) is a split sum of two
copies of \(\widehat\Z(1)\).  Here are literal quotient maps:
write \(m=Ae_{P_0}+Be_{P_1}+Ce_{P_{1/2}}\); then use
\[
 \begin{array}{c|c}
 P& M\longrightarrow M/\widehat\Z(1)e_P\simeq\widehat\Z(1)^2\\\hline
 P_0&(B,C)\\
 P_1&(A,C)\\
 P_{1/2}&(A,B)\\
 P_\infty&(A-C,B-C).
 \end{array}
\tag{A.167.12}
\]
The image of \([c]\) under the appropriate row is, by
(A.167.10), a pair of elements of \(\widehat{\Q^\times}\).
Both have zero localization in \(\widehat{\Q_p^\times}\)
at every finite prime, by the previous paragraph.  Apply exactly the
injectivity of A.76.1, taking its set of primes to be all finite
primes.  Each global completed class is zero.  This uses A.76 on its
proved inverse-limit domain, including its doubled-exponent argument;
it makes no claim of injectivity at one fixed exponent.  Since the
row (A.167.12) splits as a map of \(G_\Q\)-modules, there is a
unique class
\[
              \tau\in\widehat{\Q^\times}
                    \quad\text{with}\quad[c]=\kappa(\tau)e_P
                         \text{ in }H^1(\Q,M).
\tag{A.167.13}
\]
Uniqueness also follows from the injection of the direct summand
\(\widehat\Z(1)e_P\) followed by (A.167.9).

The local-parameter decomposition group at this rational boundary
point has the split extension
\(\widehat\Z(1)\rtimes G_\Q\), with the splitting chosen by
that parameter.  The cocycle \(\kappa(\tau)\) therefore defines
an actual full cusp section in this decomposition group.  Its image
in \(E_D/K_{\rm aff}\) is precisely (A.167.11), up to the stated
geometric conjugacy.  Conversely every such full cusp section has
actual local cusp images.  Distinct boundary points have distinct
proper Kummer classes by A.164, and (A.167.13) distinguishes their
parameters at a fixed point.  This proves the bijection (A.167.3).

\emph{The actual descent to the tripod and its tangent constants.}
The finite-cover equation and map are
\[
       aW^4=T(T-1)(T-\tfrac12),\qquad t=4T(1-T),
       \qquad W^8=\frac{t^2(1-t)}{64a^2}.
\tag{A.167.14}
\]
A section of \(E_U/K_{\rm aff}\) in the combined class
\([64a^2]\) fixes a sheet of this actual \(\mu_8\)-torsor.
Its image is consequently contained in a conjugate of
\(E_D/K_{\rm aff}\).  The local lift has the required actual
point/cusp interpretation.  Indeed a quotient conjugator lifts to
\(\Pi\), and equality of the quotient sections makes the
corresponding full local section fix that same finite sheet.  For an
ordinary point this is a rational point of its finite etale fibre.
For a cusp it is the corresponding rational boundary point with its
lifted tangent in the finite formal punctured-disc cover.  At the real
place a fixed sheet of an ordinary section is real, and a stable cusp
inertia orbit gives the real boundary section.  These are exactly the
fixed-sheet maps verified in A.157 and A.160, now before killing the
boundary inertia.  Thus (A.167.3) applies to the lifted section.

For completeness the full tangent maps, not just their inertia
exponents, are
\[
 \begin{array}{c|c|c}
 P&\text{local parameter on }D&\text{completed tangent downstairs}\\\hline
 P_0&W&\tau_t=8a\,\tau_W^4\\
 P_1&W&\tau_t=-8a\,\tau_W^4\\
 P_{1/2}&W&\tau_{1-t}=64a^2\,\tau_W^8\\
 P_\infty&W/T&\tau_{1/t}=-(a^2/4)\,\tau_{W/T}^8.
 \end{array}
\tag{A.167.15}
\]
They follow from the literal identities
\[
 \begin{aligned}
 t&=\frac{8aW^4}{1-2T}
       =\frac{-8aW^4}{2T-1},\\
 1-t&=\frac{64a^2W^8}{t^2},\\
 \frac1t&=-\frac{a^2}{4}\left(\frac WT\right)^8
       \frac1{(1-1/T)^3(1-1/(2T))^2}.
 \end{aligned}
\tag{A.167.16}
\]
In each relevant completed local ring the remaining factor has
constant term one.  For every positive integer it has the unique
root with constant term one in that characteristic-zero formal ring,
obtained recursively by division by that integer.  Its Kummer class
on the tangential section is therefore zero.  This proves all four
constants in (A.167.15), without any analytic convergence assertion.
The corresponding boundary inertia images in \(\Pi^{\rm ab}\)
are \((4,0),(4,0),(0,8),(-8,-8)\), as follows either from
(A.167.6) or from these local parameters.  The resulting combined
modulo-eight class is \([64a^2]\) in every case.  At infinity, for
example, it is represented by \(64/a^6\), whose ratio to
\(64a^2\) is the eighth power \(a^{-8}\).
Pushing the full cusp section furnished by (A.167.13) through
(A.167.14) gives the asserted actual cusp section downstairs.

\emph{The precise extra conjugacy downstairs.}
Every selected-cover section reduces to the same arithmetic section
\(E_D/H\simeq G_\Q\) of \(E_U/H\).
In this splitting the geometric group is \(C_8=\mu_8\), with its
actual cyclotomic action.  If two selected-cover sections are
conjugate by an element of \(Q_{\rm aff}\), its image
\(d\in C_8\) must centralize this arithmetic section.  Equivalently
\((\chi_8(\gamma)-1)d=0\) for every \(\gamma\in G_\Q\).
Since \(\chi_8\) has image \((\Z/8)^\times\), this gives
\[
                         d\in C_8^{G_\Q}=\{0,4\}.
\tag{A.167.17}
\]
The zero class gives exactly the conjugacy already used in
(A.167.3).  The other class is represented by the rational deck
automorphism \((T,W)\mapsto(T,-W)\).  Any base-path choice for
this automorphism changes its action by an inner conjugation from
\(H\), which disappears on the section classes under consideration.
It fixes each of the four boundary points and takes every specified
local parameter \(W\) or \(W/T\) to its negative.  Its exact
action on the completed tangent classes is therefore
\(\tau\mapsto-\tau\).  Conversely this rational deck
automorphism supplies that geometric conjugacy, so no additional
identification is missing.  The class of \(-1\) is nontrivial in
\(\widehat{\Q^\times}\), already at exponent two, and hence
each such orbit has two elements.  This proves the stated downstairs
bijection and the separation of its four boundary labels.

The integral boundary matrices, their primitive kernels, the full
semidirect-product words and all four rational tangent identities are
checked in \path{code/harmonic_affine_abelian_certificate.py}; its
executed output and exact data are retained.  The cohomological
classification is the preceding continuous argument.
\end{proof}
